\section{Related Work}
\label{sec:related_work}

\subsection{Robot Learning in Manipulation}

Robot manipulation policies have rapidly evolved from classical imitation learning methods to large-scale generalist robot policies~\citep{kim2024openvla,chen2025benchmarking}. Representative imitation learning approaches~\citep{zhao2023learning,chiDiffusionPolicyVisuomotor2024,chen2025g3flow,ze20243ddiffusionpolicygeneralizable,liu2025avr} have achieved strong visuomotor manipulation performance by learning from expert demonstrations. More recent vision-language-action models and foundation policies, including $\pi_{0.5}$~\citep{intelligence2025pi_}, RDT2~\citep{rdt2}, JoyAI-RA~\citep{zhang2026joyai}, Wall-OSS-0.5~\citep{yu2026wall}, Wall-WM~\citep{li2026wall}, LDA-1B~\citep{lyu2026lda}, Motus~\citep{bi2026motus}, G0.5~\citep{galaxea2026g05}, MotuBrain~\citep{team2026motubrain}, Qwen-RobotManip~\citep{yuan2026qwen}, Hy-Embodied-0.5-VLA~\citep{zhang2026hy} and LingBot-VA~\citep{li2026causalworldmodelingrobot}, further improve open-vocabulary instruction following, cross-task generalization, and cross-embodiment transfer. In parallel, recent studies have explored specific capability bottlenecks, such as memory-augmented manipulation~\citep{shi2025memoryvla,torne2026memmultiscaleembodiedmemory}, dynamic manipulation~\citep{cai2026internvla}, and world-model-based control~\citep{intelligence2026pi07steerablegeneralistrobotic,bi2026motus,ye2026worldactionmodelszeroshot}.

Despite this progress, existing policies still exhibit limited robustness when deployed beyond controlled task distributions. Their failures can arise from diverse factors, including insufficient temporal reasoning, weak spatial precision, limited generalization, and sensitivity to physical-world variations. This motivates evaluation protocols that go beyond aggregate task success and provide structured diagnosis across complementary manipulation capabilities and deployment settings. RoboDojo follows this direction by evaluating generalist policies through both capability-oriented simulation tasks and standardized real-world tests.

\subsection{Evaluation for Robot Manipulation}

A wide range of benchmarks have been proposed for evaluating robot manipulation. Simulation benchmarks provide scalable and reproducible testbeds: CALVIN~\citep{mees2022calvin} focuses on language-conditioned long-horizon manipulation; LIBERO~\citep{liu2023libero} studies lifelong learning across spatial, object, goal, and long-horizon suites; RoboTwin~\citep{chen2025robotwin} evaluates bimanual manipulation and generalization under structured domain randomization; RMBench~\citep{chen2026rmbench} targets memory-dependent manipulation; GarmentLab~\citep{lu2024garmentlab} focuses on garment manipulation; and SimplerEnv~\citep{li2024evaluatingrealworldrobotmanipulation} studies simulation environments that correlate with real-world performance. Real-world benchmarks, including ManipulationNet~\citep{chen2026manipulationnet}, RoboChallenge~\citep{yakefu2025robochallenge}, and RoboArena~\citep{atreya2025roboarena}, further assess policy behavior under physical deployment conditions.

These benchmarks have substantially advanced robot policy evaluation, but they are often specialized to particular task families, capability axes, or evaluation environments. Simulation-only benchmarks enable efficient and controlled stress testing, yet cannot fully capture the physical uncertainties of real-world deployment. Real-world benchmarks provide more direct evidence of deployed performance, but reproducibility remains difficult when hardware setups, scene reset procedures, evaluation protocols, and deployment interfaces vary across users or sites. RoboDojo addresses this gap by providing a unified sim-and-real benchmark with a shared policy interface and evaluation pipeline. It combines capability-oriented simulation diagnosis with reproducible real-world validation, while supporting policy integration through XPolicyLab. A detailed comparison between RoboDojo and existing benchmarks is provided in Appendix~\ref{sec:benchmark_comparison}.

% \section{Related Work}
% \label{sec:related_work}

% \subsection{Robot Learning in Manipulation}

% Robot manipulation policies have rapidly evolved from classical imitation learning methods to large-scale generalist robot policies~\citep{kim2024openvla,chen2025benchmarking}. Representative imitation learning approaches, such as ACT~\citep{zhao2023learning}, Diffusion Policy~\citep{chiDiffusionPolicyVisuomotor2024}, G3Flow~\citep{chen2025g3flow} and DP3~\citep{ze20243ddiffusionpolicygeneralizable}, have achieved strong performance on visuomotor manipulation by learning from expert demonstrations. More recent vision-language-action models and foundation policies, including $\pi_{0.5}$~\citep{intelligence2025pi_}, RDT2~\citep{rdt2}, JoyAI-RA~\citep{zhang2026joyai}, Wall-OSS-0.5~\citep{yu2026wall}, Wall-WM~\citep{li2026wall}, LDA-1B~\citep{lyu2026lda}, Qwen-RobotManip~\citep{yuan2026qwen} and LingBot-VA~\citep{li2026causalworldmodelingrobot}, further improve open-vocabulary instruction following, cross-task generalization, and cross-embodiment transfer. In parallel, recent works study specific capability bottlenecks, including memory-augmented manipulation~\citep{shi2025memoryvla,torne2026memmultiscaleembodiedmemory}, dynamic manipulation~\citep{cai2026internvla}, and world-model-based control~\citep{intelligence2026pi07steerablegeneralistrobotic,bi2026motus,ye2026worldactionmodelszeroshot}.

% Despite this progress, current robot policies remain far from human-level robustness in real-world manipulation. Many models perform well within specific task families or controlled settings, but still struggle with long-horizon execution, memory retention, fine-grained spatial reasoning, bimanual coordination, contact-rich interaction, and physical-world variations. Therefore, evaluation should not only measure task success on a narrow set of repeated patterns, but also diagnose policy reliability across complementary capability dimensions and deployment conditions. RoboDojo is designed for this purpose, combining diverse simulation tasks for rapid capability diagnosis with challenging real-world tasks for standardized physical validation.

% \subsection{Evaluation for Robot Manipulation}

% Many benchmarks have been proposed for evaluating robotic manipulation. Simulation benchmarks provide scalable and reproducible testbeds: CALVIN~\citep{mees2022calvin} focuses on language-conditioned long-horizon manipulation, LIBERO~\citep{liu2023libero} studies lifelong learning across spatial, object, goal, and long-horizon suites, RoboTwin~\citep{chen2025robotwin} evaluates bimanual manipulation and generalization under structured domain randomization, RMBench~\citep{chen2026rmbench} targets memory-dependent manipulation, GarmentLab~\citep{lu2024garmentlab} focuses on garment manipulation, and SimplerEnv~\citep{li2024evaluatingrealworldrobotmanipulation} studies simulation environments that correlate with real-world performance. Real-world benchmarks, such as ManipulationNet~\citep{chen2026manipulationnet}, RoboChallenge~\citep{yakefu2025robochallenge}, and RoboArena~\citep{atreya2025roboarena}, further evaluate policies under physical deployment conditions. However, existing benchmarks are often specialized toward a particular task family, capability axis, or evaluation environment. Simulation-only benchmarks support fast iteration and controlled stress testing, but cannot fully capture contact-rich dynamics, actuation errors, perception noise, and environmental variations. Real-world benchmarks provide more direct assessment of deployed behavior, but are costly, time-consuming, and difficult to reproduce across time, users, sites, and hardware conditions.

% RoboDojo addresses these limitations with a unified sim-and-real benchmark for generalist robot manipulation policies. It includes 42 simulation tasks and 18 real-world tasks. The simulation benchmark covers five complementary capability dimensions, including Generalization, Memory, Long-Horizon, Precision, and Open, providing diverse task designs for comprehensive diagnosis and fast model iteration. The real-world benchmark exposes policies to challenging physical-world manipulation scenarios under standardized and reproducible evaluation conditions. Through XPolicyLab, both simulation and real-world evaluation share the same policy interface and evaluation pipeline, allowing policies to be integrated once and evaluated across RoboDojo simulation and RoboDojo-RealEval with minimal policy-side adaptation. We provide a detailed comparison between RoboDojo and existing benchmarks in Appendix~\ref{sec:benchmark_comparison}.